\section{Rationale and Scope}

Bilateral stimulation is one component discussed in the EMDR literature, but the contribution of bilateral stimulation to therapeutic effects and its mechanism remain debated \parencite{landinromero2018emdr,wadji2022workingmemory}. Evidence from trauma-related dual-task paradigms cannot be transferred directly to children with ADHD, to auditory-only stimulation, or to a regulation game.

For this reason, alternating left-right audio is not part of the primary ARC-ADHD mechanism. It is a separate experimental manipulation whose value must be demonstrated against controls that distinguish bilateral alternation from ordinary sound, rhythm, expectancy, and distraction.

The paper asks one narrow question: during an otherwise identical ARC-ADHD regulation event, does alternating left-right auditory stimulation produce a specific proximal effect on recovery beyond an acoustically matched non-alternating condition?

The experiment does not test EMDR, does not involve trauma-memory processing, and does not assume that any observed effect has the same mechanism as EMDR.
